\begin{annexitem}{annex:p12.har.autarkic-eqb}{thm:p12-har-autarkic-eqb}{Every equality-code branch and the actual fixed-target typing}
\begin{proof}
The zero tester z has type nat-to-nat: its seed one has type nat,
and its step ignores both nat arguments and returns O. For the outer
eqb target \(C=\mathsf{nat}\to\mathsf{nat}\), z:C and
\(b:\mathsf{nat}\to C\to C\). Indeed, under
\(n:\mathsf{nat},q:C,m:\mathsf{nat}\), its inner recursor has seed
O:nat and step \(\lambda j:\mathsf{nat}.\lambda r:\mathsf{nat}.q j\),
of type nat-to-nat-to-nat. Application to m gives nat, and abstraction
over m gives C. This proves every formation/application premise for
eqb. The choose step similarly returns its given u:nat and ignores
the original j and recursive value r. It has target nat. Expanding
the definitions and applying the two nat constructor reductions gives
all six displayed computation equations, including the predecessor
lookup q j in the S/S equality branch.

Use ordinary nat induction on n, with the Prop motive saying that
\emph{for every m} the range disjunction and both stated equivalences
hold. These are conjunctions of propositions; no boolean value is
selected by this proof. At n=O, split the data m.
For m=O, the code is definitionally one and the equality n=m is
reflexive. Its true-code equivalence has both reflexive endpoints.
If the code is O, one=O contradicts constructor separation; this
proves the forward false-code implication by ex falso in Prop.
Conversely, \(\neg(O=O)\) contradicts reflexivity, again proving
the false-code equation by ex falso. The range is its one alternative.
For m=S j, the code computes to O and \(O=S j\) is refuted by
\(\mathsf{no\_conf}\ j\). True code would assert O=one, also
refuted by that lemma; a putative equality O=S j likewise gives FF
and hence any required true-code equation in Prop. The false-code
equivalence follows from the computed O and the proved negative
equation. The range is its O alternative.

At n=S k, first split m. At m=O the code is O and \(S k=O\)
is impossible by symmetry and \(\mathsf{no\_conf}\); the same explicit zero/one
arguments just given establish the range and both equivalences.
At m=S j, computation gives exactly \(\mathsf{eqb}\ k\ j\).
Apply the induction hypothesis at j. Successor congruence sends
k=j to S k=S j, and successor injectivity sends the latter equality
back to k=j. For the negative equivalence, from \(\neg(k=j)\)
compose with injectivity to get \(\neg(S k=S j)\); conversely
compose \(\neg(S k=S j)\) with successor congruence to get
\(\neg(k=j)\). Thus both code specifications transport to this
branch, and the range is unchanged. These four constructor cases
complete the induction. All equalities and false eliminations have
Prop targets; the decision program itself was already defined using
nat data recursion, not obtained by eliminating the proved disjunction.
\end{proof}
\end{annexitem}

\begin{annexitem}{annex:p12.har.autarkic-prime-bound}{thm:p12-har-autarkic-prime-bound}{Bounded divisor witnesses, constructed without data extraction}
\begin{proof}
Suppose one<=d. The proved additive order characterization supplies
a local Prop-existential witness k with one+k=d. The addition computation
one+k=S k, hence d=S k. The right multiplication successor equation
and commutativity of addition give
\[
 xd=x(S k)=xk+x=x+xk.
\]
The same additive order characterization with witness xk proves
x<=x+xk; transport its upper endpoint along this displayed equality
to get x<=xd. The local k existential was eliminated into this inequality
proposition, not used to produce a program parameter.
Now eliminate div d n into the stated bounded existential: for its
local coefficient x and equation xd=n, the just-proved bound transports
to x<=n. Pair that proof with the given equation and introduce the
same local x. This constructs the asserted Prop existential.

For the characterization, a forward input gives a local d with
one<d, d<n and div d n. The first strict bound implies one<=d by
the proved \(\mathsf{lthan\_leseq}\); the second implies d<=n. Apply the bounded
coefficient theorem and eliminate its Prop existential locally into
the target bounded-double-existential proposition: retain x, x<=n
and xd=n. Form the five stated conjuncts in their displayed order,
then introduce x and d. Conversely, eliminate the displayed bounded
double existential into the original proper-divisor existential.
Its local d,x and five proofs give div d n by the ex constructor
with coefficient x and equation xd=n. Pair this with the supplied
one<d and d<n to obtain propdiv d n, then introduce d.
These are all branches. The weak bounds are deliberately retained
on the right so a subsequent search may scan both d and x from O
through n inclusive and use the strict/proper conditions as data tests.
No bounded-search program or primality decision proof is inferred from
the existence of these Prop witnesses; both remain separately required.
\end{proof}
\end{annexitem}

\begin{annexitem}{annex:p12.har.autarkic-squares}{thm:p12-har-autarkic-squares}{Square growth, monotonicity and uniqueness without assuming a root}
\begin{proof}
Use only the already constructed equality, order and semiring library
of \cref{thm:p12-har-prime-equality,thm:p12-har-prime-order,thm:p12-har-prime-semiring}.
Every rewriting of an arithmetic subexpression is propositional
congruence followed by transitivity, not equality reflection.
The two multiplication successor clauses and addition associativity give
\[
 Q(S r)=r(S r)+S r=(rr+r)+S r=rr+(r+S r)=Q(r)+S(r+r).
\]
The last step is \(\mathsf{plus\_altsuc}\) at the second argument r.
For every a,u, the additive order characterization with witness u
gives \(a\le a+u\). Its successor monotonicity gives
\(S a\le S(a+u)=a+S u\). Hence \(a<a+S u\), and substituting
\(a=Q(r),u=r+r\) into the preceding equation proves square growth.

For clarity, weak/strict composition used below is itself immediate
from the actual source order: if x<=y and y<z, successor monotonicity
gives S x<=S y, then transitivity with S y<=z gives x<z.
The opposite composition, x<y and y<=z, is direct transitivity of
S x<=y<=z. Consequently neither kind of strict/weak cycle is possible,
by the proved strict irreflexivity.

Suppose a<b. Since O<=a, successor monotonicity and S a<=b imply
one<=b, i.e. O<b. Apply the additive characterization to S a<=b:
its local witness k satisfies S a+k=b, hence b=a+S k by the
addition successor clauses. Distr and commutativity give
\[
 ba=(a+S k)a=aa+(S k)a,\qquad ba=ab.
\]
Thus Q(a)<=ab, by the additive bound. Positive-factor strict-order
preservation at factor b gives ab<bb=Q(b), because O<b was proved,
not assumed. Weak/strict composition yields Q(a)<Q(b).
Every local existential coefficient in this argument is eliminated
into the final inequality proposition.

For a<=b, trichotomy gives a=b, a<b, or b<a. Equality is sent through
Q by \(\mathsf{f\_equal}\) and transports reflexivity to Q(a)<=Q(b). The second
case uses strict squaring just proved and \(\mathsf{lthan\_leseq}\). The third case
contradicts a<=b by a strict/weak cycle. This proves weak monotonicity.

Finally assume both displayed bound pairs for r,s at n. Trichotomy
of r,s has the equality branch, which is the desired conclusion.
If r<s, then S r<=s, so weak squaring gives Q(S r)<=Q(s)<=n.
Together with n<Q(S r) this is a strict/weak cycle at n. The branch
s<r is identical with the two given bound pairs exchanged:
Q(S s)<=Q(r)<=n contradicts n<Q(S s). Both impossible branches
eliminate into the equality proposition. No root was assumed to exist;
this is a uniqueness theorem for values satisfying the bounds.
\end{proof}
\end{annexitem}

\begin{annexitem}{annex:p12.har.autarkic-square-root}{thm:p12-har-autarkic-square-root}{The full F invariant and finite source-numeral computation}
\begin{proof}
Under n:nat and r:nat, Q(S r):nat; eqb(S n)(Q(S r)):nat;
and choose applied to that code, S r and r has type nat.
The displayed step therefore has exactly type nat-to-nat-to-nat,
and the nat data-recursion instance with seed O gives F:nat-to-nat.
Its constructor equations are F O=O and
\[
 F(S n)=\mathsf{choose}
   (\mathsf{eqb}\ (S n)\ (Q(S(F n))))\ (S(F n))\ (F n),
\]
where these two equations expand definitions and constructor computation.

Use ordinary nat induction into the proposition
\(P n:=(Q(F n)\le n)\wedge(n<Q(S(F n)))\).
At zero, F O=O, Q(O)=O and Q(one)=one by the actual arithmetic
computations. Reflexivity proves O<=O and one<=one is exactly
O<one. Pair them for P O.

At S n, let r abbreviate the \emph{data term} F n, put
M:=S n and q:=Q(S r), and project the induction hypothesis
into Q(r)<=n and n<q. The latter is exactly M<=q.
Use the equality-code range disjunction, eliminating into P(S n).
If eqb M q=one, its proved specification gives M=q.
Congruence for \(\lambda c:\mathsf{nat}.\mathsf{choose}\ c\ (S r)\ r\),
followed by choose's one computation, gives the \emph{propositional}
equation F(S n)=S r. For that candidate S r, its lower bound is
Q(S r)=q=M, hence reflexivity transported along this equation.
Its upper bound is M=q<Q(S(S r)), by the proved square-growth
lemma. Pair the bounds. The predicate
\(\lambda x:\mathsf{nat}.Q(x)\le M\wedge M<Q(S x)\)
is Prop-valued, so \(\mathsf{eq\_ind\_r}\) transports this pair from S r to F(S n)
along the equation F(S n)=S r. This constructs P(S n).

If eqb M q=O, its specification gives \(\neg(M=q)\).
Trichotomy of M,q, together with M<=q, now gives M<q:
equality is excluded, and q<M contradicts the weak bound by
strict/weak composition and irreflexivity. Congruence for the same
choose function, this time with its O computation, gives F(S n)=r.
The lower bound for candidate r is Q(r)<=n<=S n=M, using the
successor extension of reflexivity. Its upper bound is M<q=Q(S r),
just proved. Pair them and transport along F(S n)=r by the same
Prop-valued \(\mathsf{eq\_ind\_r}\) instance. These are both branches of the range
elimination, completing the induction.

In particular, a boolean equation was never used as definitional
conversion or as evidence from which to select data. F was defined
before either branch proof; the equations for its chosen value were
proved by \(\mathsf{f\_equal}\), and only a proposition was transported along them.
Square-root uniqueness now follows from the previous annex item.
The source expression F n+one is propositionally S(F n) by the
proved addition clause. To give its literal upper-bound formulation,
apply \(\mathsf{eq\_ind\_r}\) to the Prop predicate \(\lambda z:\mathsf{nat}.n<Q(z)\),
the obtained proof at \(z=S(F n)\), and the equation
\(F n+\mathsf{one}=S(F n)\). The output is
\(n<Q(F n+\mathsf{one})\), which pairs with the unchanged lower
bound. No propositional equation has been treated as conversion here.

For the actual numeral reduction, first note the three finite
constructor calculations. For meta-level k,l,
plus \(S^kO\) \(S^lO\) contracts to \(S^{k+l}O\), by induction
on k and the plus constructor equations. Times at those inputs
contracts to \(S^{kl}O\): at zero it is O, and its k+1 step
adds the already contracted \(S^{kl}O\) to \(S^lO\) by the first
calculation. Finally eqb at those numerals contracts to one if k=l,
and to O otherwise: repeatedly use its S/S equation until either
both arguments are O or one is O and the other a successor, then
use the corresponding explicit base equation. These are inductions
on finite literal syntax establishing reductions, not appeals to
the propositional arithmetic equalities or to an executed evaluator.

Unfold the F constructor equation from O through input45. The full
transition partition is
\[
\begin{array}{c|c|c}
 \text{next input M}&\text{current next square}&\text{root after the step}\\\hline
 1&1&1\\
 2,3&4&1\\
 4&4&2\\
 5,\ldots,8&9&2\\
 9&9&3\\
 10,\ldots,15&16&3\\
 16&16&4\\
 17,\ldots,24&25&4\\
 25&25&5\\
 26,\ldots,35&36&5\\
 36&36&6\\
 37,\ldots,45&49&6
\end{array}
\]
Initially the root is O. In each hit row the literal equality-code
calculation is one and choose returns the successor root; in each
listed miss range every literal input differs from the displayed
next square, so that calculation is O and choose returns the old root.
The literal square values in the middle column are the times
calculations for one through seven supplied above. The ranges cover
every input1 through45 without a missing transition. Thus the last
root is \(S^6O\), giving the stated finite beta/iota/delta path.
Reflexivity at that last numeral and permitted conversion give its
book-equality instance, after the universal adequacy proof has been
supplied. No actual Coq/kernel computation is reported.
\end{proof}
\end{annexitem}

\begin{annexitem}{annex:p12.har.autarkic-interpreter}{thm:p12-har-autarkic-interpreter}{The actual fold, all quotation cases and open valuation proof}
\begin{proof}
The positive constructor fields of T are: no fields for Z and U;
the nonrecursive nat field for Var; and two direct T fields for
each of Add and Mul. Accordingly the source branch generator supplies
no recursive value for Var, and exactly two, in their field order,
for each binary constructor. This constructs the five fold clauses
stated in the body. Recursion visits proper constructor subtrees only.
At any fixed target C its result is a T-to-C function; no abstraction
over an unavailable top-sort C is required. The corresponding Prop
induction receives precisely one hypothesis for each of those direct
T fields. This is the source-generated induction instance whose
branches are constructed below, not an assumption of the desired
semantic theorem or a parametricity principle.

Fix A:Set and the four specified operation/value parameters.
The E seeds are 0A:A and 1A:A. Its variable branch is
nu:nat-to-A. In a binary branch the original s,t:T and recursive
values a,b:A give a+Ab:A or a-times-Ab:A by two ordinary
applications of the supplied operation. Thus each branch really
has type T-to-T-to-A-to-A-to-A and this fold produces E-nu:T-to-A.
The same check at target T, with Var branch sigma:nat-to-T and
binary branches Add a b/Mul a b:T, types the syntactic-substitution
program. Neither fold uses a proof field or a semantic-correctness
assumption.

Prove quotation reduction by structural induction on the finite
\emph{arithmetic expression}, with its specified quotation clauses.
For source0, the quotation is Z and the E zero computation is 0A;
for source1, it is U and the E unit computation is 1A. For a source
name \(x_i\), its quotation is
\(\mathsf{Var}\ulcorner i\urcorner_{\mathsf{nat}}\), and the
variable computation is exactly
\(\nu\ulcorner i\urcorner_{\mathsf{nat}}\),
including when this is an open term in the ambient valid context.
For p+q, its quotation is Add of the two quotations. The Add
computation gives the declared addition operation applied to the
two recursively computed interpretations. Compatible reduction using
the two induction paths yields p[nu]+Aq[nu], which by the definition
of simultaneous interpretation is (p+q)[nu]. The product case is
the same actual argument with Mul and timesA. These are all five
grammar clauses; there is no omitted lambda/binder case in this
first-order arithmetic-syntax datatype. Each obtained path is finite
because the recursion descends a finite quotation tree. Its endpoint
need not be a normal form: the declared operations may be opaque.

For valuation agreement take the given pointwise proof
\(h:\Pi i:\mathsf{nat}.\nu i=\mu i\).
Apply the generated T induction to
\(P(t):=(E_\nu^A t=E_\mu^A t):\mathsf{Prop}\).
At Z and U, both interpretations compute to the same seed, so use
reflexivity on A. At Var i, both compute to the respective valuation
values and the branch is h i. At Add s t, let the two induction
proofs be e_s and e_t. Congruence for
\(\lambda z:A.z+_A(E_\nu^A t)\) applied to e_s gives
\[
 E_\nu^A s+_A E_\nu^A t=E_\mu^A s+_A E_\nu^A t.
\]
Congruence for \(\lambda z:A.(E_\mu^A s)+_A z\) applied to e_t
gives the equation from its right side to
\(E_\mu^A s+_A E_\mu^A t\). Transitivity composes them, and
the two Add computations restore exactly the branch goal P(Add s t).
For Mul, use the two lambdas with timesA instead of plusA; the same
typed \(\mathsf{f\_equal}\) and equality-transitivity terms prove P(Mul s t).
All branches conclude Prop and no equality of the functions nu,mu
themselves is needed.

To justify the finite-context valuation without postulating a global
extension, sel-A is a nat data-recursion instance at the fixed target A:
its seed v:A and step \(\lambda j:\mathsf{nat}.\lambda r:A.u\)
have the required types. Its computations are
\(\mathsf{sel}_A\ O\ u\ v\longrightarrow^*v\) and
\(\mathsf{sel}_A\ (S j)\ u\ v\longrightarrow^*u\).
Given the finite valid-context assignments, each nested application in
nu-Gamma therefore has type A, and abstraction over m:nat yields a
nat-to-A function. At literal index \(i_j\), every earlier unequal-index
eqb calculation is O by the explicit literal comparison derivation in
the F annex. These selectors contract to their remaining nesting.
The index-j calculation is one and its selector contracts to \(a_j\),
regardless of the remaining tail. Thus
\(\nu_\Gamma\ulcorner i_j\urcorner_{\mathsf{nat}}
\twoheadrightarrow_{\beta\iota\delta}a_j\).
The zero-assignment case returns the given 0A. Distinctness of the named
indices is the exact side condition used to skip earlier selectors.
This supplies a concrete extension of every finite assignment list,
not a Set-valued choice operation derived from logical existence.

For a fixed quotation using only finitely many variable indices,
the same proof recipe only calls h at a Var leaf's actual index.
Hence instead of the all-index function one may supply just the
primitive-equality proofs for that finite set of occurring indices,
and build the proof recursively: reflexivity at constants, the supplied
indexed equality at a variable, and the two explicit congruences above
at each binary node. This is a finite-context source derivation schema,
not an unimplemented runtime free-variable extractor or an assumed
functional extensionality rule. Values assigned at all other indices
are irrelevant to that quotation's derivation.

For the source's open example use the actual valuation
\(\nu_x:=\mathsf{nat\_rec}\ x\
(\lambda i:\mathsf{nat}.\lambda r:\mathsf{nat}.O)\)
in the valid context x:nat. Its zero value is x, its successor-index
values are O, and it has type nat-to-nat. Put X:=Var O.
Then the concrete quotation
\(\mathsf{Mul}(\mathsf{Add}\ X\ U)(\mathsf{Add}\ X\ U)\)
interprets, by these actual clauses, to (x+one)(x+one).
The source output quotation
\(\mathsf{Add}(\mathsf{Add}(\mathsf{Mul}\ X\ X)
(\mathsf{Mul}(\mathsf{Add}\ U\ U)X))U\)
interprets to (xx+(one+one)x)+one. This supplies the previously
omitted valuation typing and both interpretation paths, but asserts
no program G producing the second quotation from the first.
\end{proof}
\end{annexitem}

\begin{annexitem}{annex:p12.har.autarkic-substitution}{thm:p12-har-autarkic-substitution}{All five semantic-substitution branches}
\begin{proof}
Fix the actual operations, nu and sigma. The already typed E fold
gives E-nu(sigma i):A for every i:nat, so
\(\nu_\sigma:=\lambda i:\mathsf{nat}.E_\nu^A(\sigma i)\)
really has type nat-to-A. This formation is required before its use
as the right-hand interpreter's variable branch.
Use generated T induction into the precise proposition
\(P(t):=(E_\nu^A(\mathsf{sub}_\sigma t)=E_{\nu_\sigma}^A t)\).
At Z and U, syntactic substitution returns the same constructor and
both interpretations return the same seed; use reflexivity.
At Var i, substitution returns sigma i. The left side is therefore
E-nu(sigma i); the right side computes to nu-sigma i and beta-reduces
to that identical expression. Reflexivity and conversion supply the
branch, without normalizing the arbitrary sigma i syntax.

At Add s t, substitution computes to
Add(sub-sigma s)(sub-sigma t). Its E interpretation computes to
\(E_\nu^A(\mathsf{sub}_\sigma s)+_A
E_\nu^A(\mathsf{sub}_\sigma t)\).
The two induction proofs replace the first and second argument by
\(E_{\nu_\sigma}^A s\) and \(E_{\nu_\sigma}^A t\): use \(\mathsf{f\_equal}\)
for the lambda that holds the second argument fixed, then \(\mathsf{f\_equal}\)
for the lambda that holds the newly replaced first argument fixed,
and compose by equality transitivity. The result is
\(E_{\nu_\sigma}^A s+_A E_{\nu_\sigma}^A t\), precisely the
right-hand Add interpretation by its computation clause. This is the
branch goal. At Mul s t, the very same two explicitly typed congruences
with timesA instead of plusA and the Mul computation prove the goal.
These exhaust the constructors and establish the universal equation.
There is no algebraic law premise here, no semantic equality reflected
into conversion, and no function extensionality or unproved
parametricity/substitution-commutation assumption.
\end{proof}
\end{annexitem}
